\setcounter{section}{3}

  \zwischenueberschrift{Übungsaufgaben}

\inputaufgabe
{}
{

Es sei
\maabb {\gamma} { I } {\R^2
} {}
eine zweifach stetig differenzierbare
\definitionsverweis {bogenparametrisierte Kurve}{}{}
und es sei
\maabbeledisp {\delta} { -I } {\R^2
} { s } { \delta(s) = \gamma( -s)
} {,}
die umgekehrt durchlaufene Kurve. Zeige, dass die
\definitionsverweis {Krümmung}{}{}
von $\delta$ in $s$ das Negative der Krümmung von
\mathl{\gamma}{} in $-s$ ist.

}
{} {}

\inputaufgabe
{}
{

Es sei
\maabb {\gamma} { I } {\R^2
} {}
eine zweifach stetig differenzierbare
\definitionsverweis {bogenparametrisierte Kurve}{}{}
und
\maabbdisp {L} {\R^2} {\R^2
} {}
eine
\definitionsverweis {Drehung}{}{.}
Zeige, dass die
\definitionsverweis {Krümmung}{}{}
von $\gamma$ mit der Krümmung von
\mathl{L\circ \gamma}{} übereinstimmt.

}
{} {}

\inputaufgabe
{}
{

Beweise
[[Bogenparametrisierte Kurve/Ebene/Krümmungskreis/Durchlauf/Fakt|Lemma 3.6]]
im nichtstandardorientierten Fall.

}
{} {}

\inputaufgabe
{}
{

Es sei
\maabb {\gamma} { I } {\R^2
} {}
eine zweifach stetig differenzierbare
\definitionsverweis {bogenparametrisierte Kurve}{}{}
und
\maabbdisp {L} {\R^2} {\R^2
} {}
eine
\definitionsverweis {Streckung}{}{}
mit dem Streckungsfaktor

\mavergleichskette
{\vergleichskette
{ s
}
{ \neq }{ 0
}
{  }{
}
{    }{
}
{   }{
}
}
{}{}{.}
In welcher Beziehung steht die
\definitionsverweis {Krümmung}{}{}
von $\gamma$ mit der Krümmung von
\mathl{L\circ \gamma}{?}

}
{} {}

\inputaufgabe
{}
{

Es sei
\maabbeledisp {\gamma} {\R} { \R^2
} { t } { \begin{pmatrix}  \cos \left(  \omega t \right)  \\ \sin  \left(  \omega t \right)    \end{pmatrix}
} {,}
mit einem

\mavergleichskette
{\vergleichskette
{ \omega
}
{ > }{ 0
}
{  }{
}
{    }{
}
{   }{
}
}
{}{}{.}
Man zeige, dass es  unendlich viele Kreisbewegungen mit konstanter Geschwindigkeitsnorm gibt, die mit $\gamma$ für

\mavergleichskette
{\vergleichskette
{ t
}
{ = }{ 0
}
{  }{
}
{    }{
}
{   }{
}
}
{}{}{}
bis zur ersten Ableitung übereinstimmen.

}
{} {}

Die folgende Aufgabe ist wichtig für den Bau von Karussellen. Der Spaß beim Karussellfahren kommt von der Beschleunigung, nicht von der Geschwindigkeit.

\inputaufgabe
{}
{

Es sei
\maabbeledisp {\gamma} {\R} { \R^2
} { t } { \begin{pmatrix}  \cos \left(  \omega t \right)  \\ \sin  \left(  \omega t \right)    \end{pmatrix}
} {,}
mit einem

\mavergleichskette
{\vergleichskette
{ \omega
}
{ > }{ 0
}
{  }{
}
{    }{
}
{   }{
}
}
{}{}{.}
Man zeige, dass es unendlich viele Kreisbewegungen mit konstanter Geschwindigkeitsnorm gibt, die mit $\gamma$ für

\mavergleichskette
{\vergleichskette
{ t
}
{ = }{ 0
}
{  }{
}
{    }{
}
{   }{
}
}
{}{}{}
übereinstimmen und auch dort die gleiche Beschleunigung haben.

}
{} {}

\inputaufgabegibtloesung
{}
{

Es sei
\maabbeledisp {\gamma} {\R} { \R^2
} { t } { \begin{pmatrix}  \cos \left(  \omega t \right)  \\ \sin  \left(  \omega t \right)    \end{pmatrix}
} {,}
mit einem

\mavergleichskette
{\vergleichskette
{ \omega
}
{ > }{ 0
}
{  }{
}
{    }{
}
{   }{
}
}
{}{}{.}
Man zeige, dass es keine andere Kreisbewegung mit konstanter Geschwindigkeitsnorm gibt, die mit $\gamma$ für

\mavergleichskette
{\vergleichskette
{ t
}
{ = }{ 0
}
{  }{
}
{    }{
}
{   }{
}
}
{}{}{}
bis zur zweiten Ableitung übereinstimmt.

}
{} {}

\inputaufgabe
{}
{

Es sei

\mavergleichskette
{\vergleichskette
{ W
}
{ \subseteq }{ \R^2
}
{  }{
}
{    }{
}
{   }{
}
}
{}{}{}
\definitionsverweis {offen}{}{,}
\maabb {h} { W } { \R
} {}
eine zweifach
\definitionsverweis {stetig differenzierbare Funktion}{}{}
und

\mavergleichskette
{\vergleichskette
{ Y
}
{ = }{ h^{-1}(c)
}
{  }{
}
{    }{
}
{   }{
}
}
{}{}{}
die
\definitionsverweis {Faser}{}{}
zu

\mavergleichskette
{\vergleichskette
{ c
}
{ \in }{ \R
}
{  }{
}
{    }{
}
{   }{
}
}
{}{}{,}
wobei $h$ in jedem Punkt von $Y$
\definitionsverweis {regulär}{}{}
sei. Zeige, dass
\maabbdisp {\gamma} { I } { Y
} {}
genau dann eine
\definitionsverweis {geodätische Kurve}{}{}
ist, wenn $\gamma$
\definitionsverweis {bogenparametrisiert}{}{}
ist.

}
{} {}

\inputaufgabe
{}
{

\bild{ \begin{center}
\includegraphics[width=5.5cm]{ \bildeinlesungsvg {Euler spiral} {svg} }
\end{center}
\bildtext {} }

\bildlizenz { Euler spiral.svg } {} {AdiJapan} {Commons} {CC-by-sa 3.0} {}

Wir betrachten die differenzierbare Kurve
\zusatzklammer {die sogenannte \stichwort {Klothoide} {}} {} {}
\maabbeledisp {} {\R} {\R^2
} { t } { \gamma(t)
} {,}
mit

\mavergleichskettedisp
{\vergleichskette
{ \gamma(t)
}
{ =} { \sqrt{\pi} \int_0^t  \begin{pmatrix}  \cos  { \frac{   \pi s^2 }{  2 } }  \\ \sin  { \frac{   \pi s^2 }{  2 } }     \end{pmatrix}  ds
}
{ } {
}
{ } {
}
{ } {
}
}
{}{}{.}
Zeige, dass die
\definitionsverweis {Krümmung}{}{}
dieser Kurve die Gleichung

\mavergleichskettedisp
{\vergleichskette
{ \kappa(t)
}
{ =} { t
}
{ } {
}
{ } {
}
{ } {
}
}
{}{}{}
erfüllt.

}
{} {}

\inputaufgabe
{}
{

Es sei
\maabb {f} { I } {\R
} {}
eine zweifach
\definitionsverweis {stetig differenzierbare}{}{}
Funktion. Bestimme die
\definitionsverweis {Krümmung}{}{}
und den Krümmungskreis des zugehörigen Graphen.

}
{} {}

\inputaufgabe
{}
{

Es sei
\maabb {f} {\R} {\R
} {}
zweifach
\definitionsverweis {stetig differenzierbar}{}{}
und

\mavergleichskettedisp
{\vergleichskette
{ \alpha(t)
}
{ \defeq} { \begin{pmatrix}  \cos f(t)  \\ \sin f(t)   \end{pmatrix}
}
{ } {
}
{ } {
}
{ } {
}
}
{}{}{.}
Bestimme die
\definitionsverweis {Krümmung}{}{}
von $\alpha$ und den
\definitionsverweis {Krümmungskreis}{}{}
für

\mavergleichskette
{\vergleichskette
{ t
}
{ \in }{ \R
}
{  }{
}
{    }{
}
{   }{
}
}
{}{}{}
mit den Formeln aus
[[Ebene differenzierbare Kurve/Krümmung/Beschreibung/Fakt|Lemma 3.10]].

}
{} {}

\inputaufgabe
{}
{

Bestimme die
\definitionsverweis {Krümmung}{}{}
des
\definitionsverweis {Graphen}{}{}
von
\maabbdisp {\sin} {\R} {\R
} {}
für

\mavergleichskette
{\vergleichskette
{ t
}
{ = }{ { \frac{   \pi }{  2 } }
}
{  }{
}
{    }{
}
{   }{
}
}
{}{}{.}

}
{} {}

\inputaufgabe
{}
{

Bestimme die
\definitionsverweis {Krümmung}{}{}
der logarithmischen Spirale
\maabbeledisp {\gamma} {\R} { \R^2
} { t } { \gamma(t) = e^t \left(  \cos  t   , \,   \sin  t                   \right)
} {,}
für jedes

\mavergleichskette
{\vergleichskette
{ t
}
{ \in }{ \R
}
{  }{
}
{    }{
}
{   }{
}
}
{}{}{.}

}
{} {}

Eine Funktion
\maabb {f} { I } {\R
} {}
auf einem offenen Intervall

\mavergleichskette
{\vergleichskette
{ I
}
{ \subseteq }{ \R
}
{  }{
}
{    }{
}
{   }{
}
}
{}{}{}
heißt
\definitionswort {analytisch}{,}
wenn sie in jedem Punkt

\mavergleichskette
{\vergleichskette
{ a
}
{ \in }{ I
}
{  }{
}
{    }{
}
{   }{
}
}
{}{}{}
durch eine
\definitionsverweis {Potenzreihe}{}{}
beschrieben werden kann.

Eine Kurve heißt analytisch, wenn jede Komponentenfunktion analytisch ist.

\inputaufgabe
{}
{

Es sei
\maabb {\kappa} { I } {\R
} {}
eine
\definitionsverweis {analytische Funktion}{}{.}
Zeige, dass es eine analytische Kurve
\maabb {\gamma} { I } {\R^2
} {}
gibt, deren
\definitionsverweis {Krümmung}{}{}
im Punkt $\gamma(t)$ gleich $\kappa(t)$ ist.

}
{} {}

\inputaufgabe
{}
{

\bild{ \begin{center}
\includegraphics[width=5.5cm]{ \bildeinlesungsvg {Evolute-parab} {svg} }
\end{center}
\bildtext {} }

\bildlizenz { Evolute-parab.svg } {} {Ag2gaeh} {Commons} {CC-by-sa 4.0} {}

Bestimme die
\definitionsverweis {Evolute}{}{}
der Kurve
\maabbeledisp {} {\R} {\R^2
} { t } { \begin{pmatrix}  t  \\ t^2    \end{pmatrix}
} {.}

}
{} {}

\inputaufgabegibtloesung
{}
{

Bestimme die
\definitionsverweis {Krümmung}{}{}
in jedem Punkt des durch

\mavergleichskettedisp
{\vergleichskette
{ h(x,y)
}
{ =} { x^2+y^2
}
{ =} { 1
}
{ } {
}
{ } {
}
}
{}{}{}
implizit gegebenen Einheitskreises mit
[[Ebene differenzierbare Kurve/Faser/Krümmung/Fakt|Lemma 3.11]]
und dem Gradientenfeld zu $h$.

}
{} {}

  \zwischenueberschrift{Aufgaben zum Abgeben}

\inputaufgabe
{2}
{

Bestimme die
\definitionsverweis {Krümmung}{}{}
des
\definitionsverweis {Graphen}{}{}
der Exponentialfunktion
\maabbdisp {\exp} {\R} {\R
} {}
für jedes

\mavergleichskette
{\vergleichskette
{ t
}
{ \in }{ \R
}
{  }{
}
{    }{
}
{   }{
}
}
{}{}{.}

}
{} {}

\inputaufgabe
{2}
{

Bestimme die
\definitionsverweis {Krümmung}{}{}
der archimedischen Spirale
\maabbeledisp {\gamma} {\R} { \R^2
} { t } { \gamma(t) = t \left(  \cos  t   , \,   \sin  t                   \right)
} {,}
für jedes

\mavergleichskette
{\vergleichskette
{ t
}
{ \in }{ \R
}
{  }{
}
{    }{
}
{   }{
}
}
{}{}{.}

}
{} {}

\inputaufgabe
{4}
{

Man gebe ein Beispiel für eine zweifach stetig differenzierbare bogenparametrisierte Kurve
\maabbdisp {\gamma} {\R} {\R^2
} {}
derart, dass für zwei Zeitpunkte

\mavergleichskette
{\vergleichskette
{ t_0
}
{ \neq }{ t_1
}
{  }{
}
{    }{
}
{   }{
}
}
{}{}{}
die Gleichheit

\mavergleichskette
{\vergleichskette
{ \gamma (t_0)
}
{ = }{ \gamma(t_1)
}
{  }{
}
{    }{
}
{   }{
}
}
{}{}{}
gilt und derart, dass die
\definitionsverweis {Krümmungskreise}{}{}
zu
\mathkor {} {t_0} {und} {t_1} {}
nicht übereinstimmen.

}
{} {}

\inputaufgabe
{4}
{

Bestimme die
\definitionsverweis {Evolute}{}{}
der Kurve
\maabbeledisp {} {\R} {\R^2
} { t } { \begin{pmatrix}  2 \cos  t  \\ 3 \sin  t    \end{pmatrix}
} {.}
In welchen Punkten hat die Ableitung der Evolute Nullstellen?

}
{} {}

\inputaufgabe
{4}
{

Bestimme die
\definitionsverweis {Krümmung}{}{}
des durch

\mavergleichskettedisp
{\vergleichskette
{ h(x,y)
}
{ =} { x^4+y^4
}
{ =} { 1
}
{ } {
}
{ } {
}
}
{}{}{}
implizit gegebenen Einheitskreises mit
[[Ebene differenzierbare Kurve/Faser/Krümmung/Fakt|Lemma 3.11]]
und dem Gradientenfeld zu $h$ in den folgenden Punkten.
\aufzaehlungzwei {
\mavergleichskette
{\vergleichskette
{ P
}
{ = }{ \begin{pmatrix}  1  \\ 0   \end{pmatrix}
}
{  }{
}
{    }{
}
{   }{
}
}
{}{}{,}
} {
\mavergleichskette
{\vergleichskette
{ Q
}
{ = }{ \begin{pmatrix}  { \frac{   1  }{  \sqrt[4]{2}  } }  \\ { \frac{   1  }{  \sqrt[4]{2}  } }     \end{pmatrix}
}
{  }{
}
{    }{
}
{   }{
}
}
{}{}{.}
}

}
{} {}

 [[Kategorie:Latexseite]]
